\section{Background}

\subsection{FEMA National Flood Hazard Layer}

The FEMA National Flood Hazard Layer (NFHL) is the authoritative federal dataset defining
flood zone designations across the United States \citep{FEMA_NFHL}. For the purposes of
highway vulnerability assessment, the critical designation is the Special Flood Hazard Area
(SFHA) --- the land subject to inundation by a flood with a 1\% annual probability of
occurrence (the ``100-year flood''). SFHA zones include Zone A (riverine flood), Zone AE
(detailed study areas with base flood elevations established), Zone AH (shallow flooding),
Zone AO (sheet flow), and Zone VE (coastal areas with wave action). For corridor scoring,
we treat all SFHA zone types equivalently, as each represents a 1\% annual inundation
probability that would render the highway surface inaccessible.

The NFHL is delivered via FEMA's Map Service Center API, which supports polygon overlay
queries against arbitrary geographic buffers. We construct a 1-mile buffer on each side of
each interstate centerline (2-mile total corridor envelope) and intersect against NFHL SFHA
polygons to compute SFHA coverage in corridor miles. This buffer captures not just the
roadway surface but also the drainage infrastructure and access roads that must remain
operational for the corridor to function.

A critical limitation of static NFHL data: maps reflect conditions at their last revision
date, which varies significantly by county. In rapidly subsiding areas (coastal Louisiana,
Houston subsidence zone), the effective flood zone is larger than the NFHL polygon suggests,
because land has subsided since the last map revision. We address this through 2050 projection
methodology (Section~\ref{sec:projections}).

\subsection{NOAA Sea Level Rise Projections}

NOAA's 2022 Sea Level Rise Technical Report provides regionally resolved SLR projections
through 2100 under five scenarios \citep{NOAA_SLR2022}. For this analysis, we use two
scenarios: the Intermediate scenario (0.5m global mean SLR by 2050, roughly corresponding
to SSP2-4.5) and the High scenario (1.0m by 2050, corresponding to SSP5-8.5). NOAA's
regional adjustments account for ocean dynamics, gravitational fingerprinting, and local
land subsidence, producing spatially heterogeneous projections even under a single global
scenario. Gulf Coast regions under the Intermediate scenario show 0.6--0.7m effective local
SLR by 2050 due to additional land subsidence; the 0.5m global figure understates Gulf Coast
exposure.

These projections are used here to project SFHA expansion: we apply the projected SLR to
current elevation data (1/3 arc-second National Elevation Dataset) and recompute SFHA
polygon boundaries at projected water levels, producing 2050 corridor vulnerability scores.

\subsection{Historical Closure Data}

For winter precipitation and wildfire corridors, exposure cannot be captured by static
flood-zone overlay. We use three sources of historical closure frequency data:

\textbf{Caltrans District incident database} \citep{Caltrans2023}: California corridor
closures, 2010--2022. Includes I-80 Donner Pass, I-5 Siskiyou Summit, US-50 Echo Summit.
Reported in closure events per year with duration in hours. Donner Pass (I-80): 50 closures
per year mean, 18 hour mean duration, maximum recorded closure 72+ hours (January 2019 storm).

\textbf{WSDOT Incident Database}: Washington state corridors, 2012--2022. I-90 Snoqualmie
Pass: 40 closures per year mean, 12 hour mean duration.

\textbf{FHWA National Incident Database} \citep{FHWA_incident2022}: national coverage for
major incidents (duration $\geq$ 2 hours). Used for I-35 Oklahoma tornado corridor, I-70
Vail Pass, I-90 Homestake (Montana).

\subsection{Climate--Transportation Literature}

The relationship between climate exposure and transportation network performance has received
growing attention since the 2005 Katrina disruption of Gulf Coast I-10, which caused an
estimated \$1.5 billion in freight disruption costs \citep{FHWA_incident2022}.
\citet{pregnolato2017impact} developed the first comprehensive framework for quantifying
flood impacts on transport network accessibility, showing that accessibility loss is a
superlinear function of inundation depth on critical network links. The FHWA vulnerability
assessment framework \citep{FHWA_vulnerability2012} established a four-step methodology
for identifying climate-vulnerable highway assets, but did not produce corridor-level
exposure scores or investment prioritization guidance. The National Academies report on
extreme weather and the transportation system \citep{NationalAcademies2016} identified sea
level rise, extreme precipitation, and wildfire smoke as the three primary climate stressors
for the interstate highway system, and called for standardized exposure metrics --- a call
that remains unmet in federal investment criteria as of 2026.

\subsection{Two Exposure Mechanisms: Flood vs. Winter/Fire}

The D1 dimension synthesizes two structurally different exposure mechanisms. Flood exposure
(Gulf Coast, Atlantic coastal corridors) is structural: the road is built in a location that
is, with some annual probability, inundated. The metric is geometric --- how many miles of
road lie within the flood plain. Winter precipitation and wildfire exposure, by contrast,
is operational: the road is geographically above the flood plain, but weather conditions
render it intermittently impassable. The metric is frequency-duration: how often does the
road close and for how long.

Combining these two mechanisms into a single D1 score requires a calibration that equates
structural flood exposure with operational weather exposure. We anchor the calibration on
Donner Pass: 50 closures per year averaging 18 hours each (900 annual closure-hours) maps
to D1 = 7.8 under the operational exposure branch of the scoring function. Gulf Coast
I-10's consecutive-SFHA overlay produces D1 = 8.4 under the structural exposure branch.
This calibration implies that 127 consecutive SFHA miles is somewhat worse than 900 annual
closure-hours --- a judgment anchored in the observation that flood exposure eliminates
the corridor entirely during events that may recur annually, while Donner closures, though
frequent, are recoverable within hours to days.
